\subsection{How to read the row expansion}
The weight of a square depends on how it enters the mean in the preceding table. For terms that
average over a center \(c\), we use the measure \(\nu(c)\,dc\) and write \(\langle\cdot\rangle_c\). A
second independent square \(i\) carries the same measure. If \(i\) lies on a line \(l\) through the
center, only one coordinate is free, either row or column, and the measure on that fiber is
\(d\omega_l(i)=\eta_l\nu(i)\,dy\). In this paragraph \(C\) denotes a kernel distinct from the column
family. Define
\begin{align*}
 A_{lh}&=(L_0-K)_{lh},\qquad B_{lh}=K_{lh}(W_{lh}-K_{lh}),\\
 C_{lh}&=-B_{lh}(W_{lh}-2K_{lh}),\qquad
 w(i)=\sum_{\{l,h\}\subset J(i),\,l\ne h}(L_0)_{lh}.
\end{align*}
Here \(W_{lh}\) equals \(W\) at the crossing, or zero if there is no crossing. We write
\(A(i,c)=\sum_{l\in J(i),\,h\in J(c)}A_{lh}\) and similarly for \(B,C\). If the first argument is a
single line, we sum only over \(h\). After summing, we set \(M_2=A^2+B,\ M_3=A^3+3AB+C\). The cubic
term distinguishes three averages according to whether we choose an independent square, a diagonal
line or a square on a fiber through the center. With their respective signs we obtain
\begin{align*}
 f2E&=\left\langle -M_3(i,c)/6+\tfrac32 W(i)M_2(i,c)\right\rangle_{i,c},\\
 f2z&=-\tfrac16\left\langle
 \int_{l\in D\cup A}\frac{u_l}{U_l}M_3(l,c)\,d\rho_l\right\rangle_c ,\\
 f2F&=\left\langle\sum_{l\in J(c)}\int_{i\in l}W(i)^2\,d\omega_l(i)\right\rangle_c .
\end{align*}
The measure \(d\rho\) is the continuum measure represented by the line weights in
Section~\ref{sec:profile-data}. Four other terms use the same averages.
\begin{align*}
 G_{\rm cof}&=-\tfrac12-3\langle W^2\rangle_c,\qquad
 f4_{\rm expl}=-9\langle W^2\rangle_c,\\
 f4_{L0}&=\langle Ww\rangle_c,\qquad f5c=-3\langle Ww\rangle_c .
\end{align*}
For \(C3_{\rm cuad}\) we consider a center and three translated squares that complete the geometric
square. We take both oriented bases (row--column and diagonal--antidiagonal) and both signs in each
direction. The weight is the product of the four \(W\) factors, \(x^{n_{RC}-3}\) and the \(U_l\)
factors of the distinct diagonal lines. Here \(n_{RC}\) counts the distinct rows and columns. We
integrate with respect to Lebesgue measure \(dr\,dc\,dh\), with \(h>0\) and the square inside the
board. The four \(W\) factors already supply the physical weight, with no additional density at the
center. In the grid implementation a side of length zero carries half weight. In \(f3F\) the quadratic
fiber term includes the variance and the covariance due to shared lines, and \(f3c\) collects the two
center contractions. In \(f5F\) we integrate over the fiber the product of \(-w(i)\) and the sum of
its line responses. In \(f7\) we square \emph{after} summing the responses on the fiber and outside
the center. Finally, \(Q1,Q2\) take the linear and quadratic trace terms with respect to the same
center, with the corrections \(\Delta_1,\Delta_2\) to the product over distinct lines given above.
These measures and multiplicities determine the interior integrand \(d(x)\). At arrival we must
average over the diagonal pairs that remain available in the rook mask.
